\documentclass[11pt]{article}

\usepackage[T1]{fontenc}
\usepackage[margin=1.12in]{geometry}
\usepackage{amsmath,amssymb,amsthm,mathtools}
\usepackage{xcolor}
\usepackage[colorlinks=true,linkcolor=blue,citecolor=teal,urlcolor=magenta]{hyperref}
\usepackage[nameinlink,capitalise]{cleveref}

\newtheorem{proposition}{Proposition}
\newtheorem{remark}{Remark}

\newcommand{\A}{\mathbb A}
\newcommand{\Pj}{\mathbb P}
\newcommand{\F}{\mathcal F}
\newcommand{\cP}{\mathcal P}
\newcommand{\m}{\mathfrak m}
\newcommand{\Pic}{\operatorname{Pic}}
\newcommand{\Spec}{\operatorname{Spec}}
\newcommand{\Fr}{\operatorname{Fr}}

\title{A Tendler-Style Proof of Extension\\Across Takeuchi's Boundary}
\author{Standalone note}
\date{July 2026}

\begin{document}

\maketitle

\begin{abstract}
This note isolates a short proof that a rank-one local system which is tame
along the boundary of Takeuchi's compactified Abel--Jacobi morphism extends
across that boundary.  The proof adapts the power argument in
Tendler's treatment of tame geometric class field theory: a prime-to-the-
characteristic power kills tame inertia, while the abelian fundamental group
of affine space in positive characteristic is pro-$p$.  The argument applies
to a nonreduced modulus only after the ramification calculation in Chapter~2
has made the symmetric-power local system tame at the full boundary.
\end{abstract}

\section{Setting}

Let $k$ be a perfect field, let $C$ be a projective smooth geometrically
connected curve over $k$, let $\m>0$ be an effective divisor, and put
$U=C\setminus |\m|$.  Let $\Lambda$ be a finite ring whose cardinality is
invertible in $k$, and let $\F$ be a locally constant sheaf of free rank-one
$\Lambda$-modules on $U$.

For sufficiently large $d$, write
\[
 \Phi_d\colon U^{(d)}\longrightarrow \Pic_{C,\m}^{d}
\]
for the Abel--Jacobi morphism and
\[
 \widetilde\Phi_d\colon
 \widetilde C^{(d)}_{\m}\longrightarrow \Pic_{C,\m}^{d}
\]
for Takeuchi's compactification.  It is a projective-space bundle.  If
\[
 H=\widetilde C^{(d)}_{\m}\setminus U^{(d)},
\]
then, over a geometric generic point of $\Pic_{C,\m}^{d}$, the pair consisting
of the compactified fiber and its boundary is
\begin{equation}\label{eq:generic-fiber}
 \left(\Pj^r,\Pj^{r-1}\right),
 \qquad
 r=d-\deg(\m)-g+1,
\end{equation}
and its open fiber is $\A^r$.

The input supplied by the ramification calculation is:
\begin{equation}\label{eq:tame-input}
 \F^{[d]}\text{ is tamely ramified at the generic point of }H.
\end{equation}

\section{The extension argument}

\begin{proposition}\label{prop:extension}
Assume \eqref{eq:tame-input}.  Then $\F^{[d]}$ extends uniquely to a
locally constant sheaf of free rank-one $\Lambda$-modules on
$\widetilde C^{(d)}_{\m}$.
\end{proposition}

\begin{proof}
Let $\xi$ be the generic point of $\Pic_{C,\m}^{d}$ and set
$K=\overline{k(\xi)}$.  By \eqref{eq:generic-fiber}, the geometric generic
open fiber is
\[
 V=\A^r_K\subset \overline V=\Pj^r_K,
 \qquad
 D_\infty=\overline V\setminus V=\Pj^{r-1}_K.
\]
Let $\cP^{[d]}=\Fr(\F^{[d]})$ be the $\Lambda^\times$-torsor of frames.
Its restriction to $V$ corresponds to a continuous character
\[
 \chi\colon \pi_1(V)\longrightarrow \Lambda^\times.
\]
Because $\Lambda^\times$ is finite, the image of $\chi$ is finite.

Suppose first that $\operatorname{char}k=0$.  Then $\pi_1(\A^r_K)=1$;
see \cite[Example~4.14]{Tendler}.  Hence $\chi$ is trivial.

Now suppose that $\operatorname{char}k=p>0$.  Let $I_\infty$ denote the
inertia group at the generic point of $D_\infty$.  By the tameness assumption,
the finite group $\chi(I_\infty)$ has order prime to $p$.  Let $N$ be its
exponent.  Then
\[
 \gcd(N,p)=1
 \qquad\text{and}\qquad
 \chi^N|_{I_\infty}=1.
\]
Thus the rank-one local system corresponding to $\chi^N$ is unramified at
the generic point of $D_\infty$.  Since $\Pj^r_K$ is regular and
$D_\infty$ is the only codimension-one component of its complement, purity
of the branch locus extends this local system from $V$ to $\Pj^r_K$.
Projective space is simply connected, so the extension is constant.  Hence
\begin{equation}\label{eq:N-power-trivial}
 \chi^N=1\quad\text{on }\pi_1(V).
\end{equation}

Since $\chi$ is rank one, it factors through $\pi_1(V)^{\mathrm{ab}}$.
Tendler proves that
\[
 \pi_1(\A^r_K)^{\mathrm{ab}}
 \quad\text{is a pro-}p\text{ group}
\]
in \cite[Example~4.15]{Tendler}.  Consequently, the finite group
$\operatorname{im}(\chi)$ is a $p$-group.  On the other hand,
\eqref{eq:N-power-trivial} says that its exponent divides $N$, which is
prime to $p$.  The only such $p$-group is the trivial group.  Therefore
$\chi=1$.

We have shown in either characteristic that $\cP^{[d]}$ is constant on the
geometric generic affine fiber.  In particular, its inertia at the generic
point of the hyperplane at infinity is trivial.  Unramifiedness can be
checked after the separable extension $k(\xi)\subset K$, so
$\cP^{[d]}$ is unramified at the generic point of $H$.

Finally, $\widetilde C^{(d)}_{\m}$ is regular and $H$ is the only
codimension-one component of its complement $U^{(d)}$.  Purity of the
branch locus therefore extends $\cP^{[d]}$ to a finite etale
$\Lambda^\times$-torsor on $\widetilde C^{(d)}_{\m}$.  The associated
rank-one $\Lambda$-module is the required extension of $\F^{[d]}$.
Uniqueness follows because a dense open immersion into a normal connected
scheme induces a surjection on etale fundamental groups.
\end{proof}

\section{Relation with Tendler's proof}

In Tendler's tame argument, a character $\rho$ is replaced by a power
$\rho^{q-1}$ which is unramified.  The unramified geometric class field
theory theorem makes the induced power of the symmetric-product character
trivial on every affine Abel--Jacobi fiber.  Tendler then uses the fact that
$\pi_1(\A^r)^{\mathrm{ab}}$ is pro-$p$ and that $q-1$ is prime to $p$ to
deduce that the original character is already trivial on the fiber; see
\cite[\S9.3]{Tendler}.

The argument above makes two adjustments.

\begin{enumerate}
 \item The base field here is an arbitrary perfect field, so there is no
 distinguished integer $q-1$.  Instead, $N$ is chosen to be the exponent of
 the actual finite tame-inertia image.  Finiteness of $\Lambda$ guarantees
 that such an $N$ exists, and tameness guarantees that $N$ is prime to $p$.

 \item The original local system $\F$ may have wild ramification allowed by
 the nonreduced modulus $\m$.  Therefore no prime-to-$p$ power of $\F$ need
 be unramified on $C$.  The power argument is applied only after passing to
 $\F^{[d]}$ and to the full boundary $H$, where the Chapter~2 calculation
 has already proved tameness.
\end{enumerate}

Thus Tendler's method does not replace the ramification calculation.  It
replaces the final step from ``tame at the full boundary'' to ``unramified,
hence extendable across the full boundary.''

\begin{remark}
For rank-one local systems, this power argument proves exactly the special
case needed here of the statement that affine space has trivial tame
fundamental group relative to the hyperplane at infinity.  It therefore
provides a more explicit alternative to citing that statement directly.
\end{remark}

\begin{thebibliography}{9}

\bibitem{Tendler}
Avichai Tendler,
\emph{Geometric Class Field Theory},
Master's thesis, Hebrew University of Jerusalem, 2010;
arXiv:1507.00104 (2015).
\url{https://arxiv.org/abs/1507.00104}.

\bibitem{SGA1}
Alexander Grothendieck,
\emph{Rev\^etements \'etales et groupe fondamental (SGA 1)},
Lecture Notes in Mathematics 224, Springer, 1971.

\end{thebibliography}

\end{document}
